\documentclass[letterpaper,10pt]{article}
\usepackage[utf8]{inputenc}
\usepackage[T1]{fontenc}
\usepackage{lmodern}
\usepackage{microtype}
\usepackage[margin=0.45in]{geometry}
\usepackage{hyperref}
\usepackage{titlesec}
\usepackage{enumitem}
\usepackage{xcolor}

\hypersetup{
    colorlinks=true,
    linkcolor=blue,
    urlcolor=blue,
    pdftitle={Saul Cordoba - CV Data BI},
    pdfauthor={Saul Cordoba},
    pdfsubject={Data Engineer \& BI Analyst | Tableau, Power BI, SQL, Python},
}

\pagestyle{empty}
\urlstyle{same}

\titleformat{\section}{\large\bfseries}{}{0em}{}[\titlerule]
\titlespacing*{\section}{0pt}{5pt}{3pt}

\setlist[itemize]{leftmargin=1.2em, itemsep=1pt, topsep=1.5pt, parsep=0pt, partopsep=0pt}

\begin{document}

%--- HEADER ---
\begin{center}
    {\Huge \bfseries SAÚL CÓRDOBA} \\[3pt]
    {\Large \textbf{Data Engineer \& BI Analyst | Tableau, Power BI, SQL, Python}} \\[2pt]
    Villavicencio, Colombia (Disponible para Trabajo Remoto / Relocalización) \\
    (+57) 321 961 1182 | \href{mailto:sauleduardo1111@gmail.com}{sauleduardo1111@gmail.com} | \href{https://github.com/Garetetete}{GitHub: Garetetete} | \href{https://www.linkedin.com/in/saul-cordoba-a2997427b}{LinkedIn} | \textbf{Inglés:} B2
\end{center}

%--- PERFIL PROFESIONAL ---
\section{PERFIL PROFESIONAL}
\textbf{Ingeniero de datos y analista BI} con experiencia en el \textbf{sector petrolero de EE.UU.} — consultas \textbf{SQL sobre Oracle} y dashboards en \textbf{Tableau y Power BI} — y en pipelines de datos a escala: ingesta desde \textbf{APIs de datos abiertos gubernamentales} (Socrata) hacia \textbf{PostgreSQL} con visualización en \textbf{Looker Studio} y KPIs de negocio. Fuerte base de ingeniería de software (Python, SQL, Docker) para construir pipelines reproducibles de extremo a extremo. Ingeniero Electrónico con doble especialización (Instrumentación y Control Industrial; Ingeniería de Software).

%--- HABILIDADES TÉCNICAS ---
\section{HABILIDADES TÉCNICAS}
\begin{itemize}
    \item \textbf{BI y visualización:} Tableau, Power BI, Looker Studio, Chart.js, diseño de KPIs y metas de negocio.
    \item \textbf{Bases de datos y SQL:} \textbf{Oracle}, PostgreSQL (JSONB, índices, migraciones, Row-Level Security), SQL Server, MongoDB, Redis, consultas analíticas y modelado de datos.
    \item \textbf{Ingeniería de datos (Python):} Pandas, openpyxl (automatización de Excel), SQLAlchemy, limpieza y normalización de datos (UTF-8/acentos), scraping con Selenium, APIs REST y Socrata/SoQL (paginación de 50k filas, particionado temporal).
    \item \textbf{IA aplicada a datos:} clasificación y extracción de información con LLMs (Vertex AI/Gemini), salida estructurada JSON, validación cruzada con modelos secundarios.
    \item \textbf{Infraestructura:} Docker y Docker Compose, Linux, ngrok, Azure DevOps, GitHub Actions, CI/CD.
\end{itemize}

%--- EXPERIENCIA ---
\section{EXPERIENCIA PROFESIONAL}

\textbf{Desarrollador de Software} \hfill \textbf{2025 -- Presente} \\
\textit{SASOFTCO} \hfill \textit{Villavicencio, Colombia}
\begin{itemize}
    \item \textbf{Pipeline de datos de una plataforma SaaS en producción}, que \textbf{creé desde cero de extremo a extremo}: ingesta automatizada desde \textbf{APIs de datos abiertos gubernamentales} (Socrata) — consultas SoQL paginadas a 50.000 filas, particionado por rango temporal para evitar timeouts y filtrado por múltiples criterios — más APIs comerciales, hacia PostgreSQL 15 como fuente única de verdad.
    \item \textbf{Enriquecimiento con IA}: clasificación y extracción estructurada de documentos con Vertex AI/Gemini, series históricas de precios y consolidación de datos para comparación automatizada.
    \item \textbf{Dashboards de KPIs y metas anuales} con visualizaciones Chart.js para seguimiento comercial; import/export automatizado de Excel (openpyxl) para integración con flujos legacy.
    \item \textbf{Pipeline CSV $\rightarrow$ PostgreSQL $\rightarrow$ Looker Studio}: limpieza y normalización de datos (UTF-8/acentos), carga a PostgreSQL dockerizado y exposición a Looker Studio.
    \item \textbf{Extracción documental automatizada}: Python + Selenium + Pandas sobre portales web, con mapeo inteligente de campos, descarga masiva de CSV/PDF y filtrado con IA.
\end{itemize}

\textbf{Desarrollador / Analista de Datos} \hfill \textbf{2024 -- 2025} \\
\textit{Btilink SAS} \hfill \textit{Villavicencio, Colombia}
\begin{itemize}
    \item \textbf{Analista de datos para Aera Energy LLC} (compañía petrolera de California, EE.UU.): consultas \textbf{SQL sobre bases de datos Oracle} para extraer y consolidar información de la operación petrolera.
    \item Dashboards y reportes en \textbf{Tableau y Power BI} sobre volúmenes de producción, ingresos y salidas de la operación.
    \item Reportes utilizados por múltiples equipos para el \textbf{monitoreo continuo del negocio}.
\end{itemize}

%--- EDUCACIÓN ---
\section{EDUCACIÓN}
\textbf{Especialización en Ingeniería de Software} \hfill \textbf{2025 -- 2026} \\
Universidad de los Llanos \\
\textbf{Especialización en Instrumentación y Control Industrial} \hfill \textbf{2024 -- 2025} \\
Universidad de los Llanos --- \textit{Modalidad de grado: Diplomado en Sistemas IoT y PLC} \\
\textbf{Ingeniería Electrónica} \hfill \textbf{2018 -- 2024} \\
Universidad de los Llanos

%--- CERTIFICACIONES (descomentar y completar cuando tengas los datos) ---
% \section{CERTIFICACIONES Y DISTINCIONES}
% \begin{itemize}
%     \item \textbf{Nombre de la certificación} -- Entidad emisora (Año).
%     \item \textbf{Distinción académica:} promedio destacado, beca o reconocimiento.
% \end{itemize}

\end{document}
